\documentclass[12pt]{article}
\usepackage[utf8]{inputenc}
\usepackage[T1]{fontenc}
\usepackage[margin=1in]{geometry}
\usepackage{setspace}
\usepackage{xcolor}
\usepackage[colorlinks=true, urlcolor=blue!60!black, linkcolor=blue!60!black]{hyperref}
\setlength{\parindent}{0em}
\setlength{\parskip}{0.6em}
\renewcommand{\baselinestretch}{1.0}
\begin{document}
\begin{center}
{\LARGE \textit{Katzenbach v. Morgan}, 384 U.S. 641 (1966)\par}
\vspace{0.5em}
{\large Case Brief --- Two-Page Edition\par}
\end{center}

\section{Memory Jogger}
Section 5 of the Fourteenth Amendment is a positive grant of power under the \textit{McCulloch} standard, so Congress could displace New York's English literacy test for Puerto Rican voters without a prior judicial ruling against it (\href{https://supreme.justia.com/cases/federal/us/384/641/}{Justia}).

\section{Detailed Case Facts}
New York required new voters to read and write English, disenfranchising many Puerto Rican migrants schooled in Spanish. Section 4(e) of the Voting Rights Act of 1965 barred denying the vote to anyone with a sixth-grade education in an ``American-flag'' school in Puerto Rico. \textit{Lassiter} (1959) had upheld literacy tests against judicial challenge. New York City voters sued, arguing \S 4(e) exceeded Congress's power.

\section{Procedural History}
A three-judge District Court for the District of Columbia held 2--1 that Congress exceeded its powers and enjoined \S 4(e). The Attorney General and the city Board of Elections appealed directly.

\section{Judicial Votes}
Reversed, 7--2. Brennan for the Court (Warren, Black, Clark, White, Fortas; Douglas in part); Harlan, with Stewart, dissenting.

\section{Holding}
Section 4(e) is appropriate \S 5 legislation and preempts New York's test. Congress need not await a judicial finding; it suffices that the Court can perceive a rational basis for Congress's judgment. \textit{City of Boerne} (1997) later narrowed \textit{Morgan} to its remedial rationale.

\section{Analysis of Opinions}
\textbf{Brennan:} A judicial-finding requirement would reduce Congress to ``the insignificant role of abrogating only those state laws that the judicial branch was prepared to adjudge unconstitutional.'' Congress could find the vote would protect Puerto Ricans in public services, or could itself judge the test invidious. Footnote 10: Congress may expand, not dilute, rights. \textbf{Harlan:} whether a state law infringes the Constitution is ``a judicial question''; Congress built no record and could equally dilute rights.

\section{Examples: Future Applications}
\textbf{Same-Side:} A nationwide literacy-test ban on a record of segregated schooling (\textit{Oregon v. Mitchell}). \textbf{Opposite-Side:} RFRA's override of a Free Exercise ruling, struck down in \textit{Boerne}. \textbf{Fence-Sitter:} Lowering the voting age to 18 in state elections, which split the Court 5--4.

\section{Critique}
Critics say the power to ``\textit{enforce}'' presupposes a fixed baseline and that the one-way ratchet is unworkable. Post and Siegel defend \textit{Morgan} as democratic constitutionalism against ``juricentric'' limits.

\section{Key Quotations}
\begin{quote}
``Thus, the \textit{McCulloch v. Maryland} standard is the measure of what constitutes `appropriate legislation' under \S 5 of the Fourteenth Amendment. Correctly viewed, \S 5 is a positive grant of legislative power authorizing Congress to exercise its discretion in determining whether and what legislation is needed to secure the guarantees of the Fourteenth Amendment.''
\end{quote}
\noindent --- Justice Brennan, Opinion of the Court, 384 U.S. at 651.

\end{document}
